\section{Operační systém}

Pro zajištění co největší efektivity je důležité si vybrat správný operační systém. Samozřejmě se nabízí možnost použít Windows, který berou vývojáři assetů a doplňků nejvíce v potaz, jelikož je to stále nejpoužívanější systém pro běžné uživatele. Ovšem u tohoto systému je obrovský problém s optimalizací a efektivitou. Windows má otevřenou spoustu aplikací na pozadí, které značně brzdí celkový výkon počítače. Z těchto důvodů byla v našem týmu provedena migrace na operační systém Linux. V našem týmu jsou aktuálně využívány distribuce Arch Linux, Fedora Linux a Ubuntu. Rychlost a efektivita v těchto Linux distribucích je mnohonásobně vyšší než u Windows.

\section{Herní engine}

Herní engine dává vývojářům framework pro vývoj her bez nutnosti vytváření všech systémů, jako například fyziku, grafiku nebo AI od úplného začátku [1]. Herní engine odstraňuje potřebu programovat a integrovat systémy pro podporu hry, tímto šetří čas a zdroje vývojářů, aby se mohli soustředit na tvorbu samotné hry. Na trhu je několik různých možností enginů na výběr. Mezi nejznámější patří třeba Unreal Engine 5, Unity Engine nebo Godot. Tento projekt používá Unity Engine z několika důvodů. Unity je stále jeden z nejpopulárnějších enginů současnosti, díky tomu existuje spousta materiálů, které může člověk při vývoji použít. Unity je zároveň vhodné jak pro tvorbu 3D, tak pro 2D projektů, na rozdíl od Unreal Engine 5, který se specializuje především na 3D. Další důležité kritérium při výběru herního enginu je jazyk, který se používá pro programování. Unreal používá C++, Godot má možnost použít jejich vlastní GDScript, který se nejvíce podobá jazyku Python, ale nabízí také možnost použít C\#, C nebo C++ díky GDExtension. Unity používá defaultně C\#, což je zároveň jazyk, který se, jakožto studenti IT na SPŠ Karviná, učíme od prvního ročníku v předmětu \texttt{programování}. Bohužel v době vybírání enginu ještě Godot nebyl ani zdaleka tak pokročilý, jako je v dnešní době. Kdyby byl Godot už tehdy stabilní, tak by nejspíše byl využit i pro tento projekt. Ovšem díky době, komunitě, jazyku a možnostem byl pro tento projekt zvolen engine Unity.

\section{Programovací jazyk}

V herních enginech se programuje hlavně pomocí programovacích jazyků [2]. Avšak třeba Unreal nabízí možnost programovat čistě pomocí bloků, kterými se nastavuje chování a spojují se mezi sebou. Vzhledem k tomu, že tento projekt jede na Unity, tak používá jazyk C\#. C\# je moderní a intuitivní programovací jazyk, který vývojářům umožňuje vytvářet mnoho typů bezpečných a stabilních aplikací. Jeho kořeny sahají do rodiny jazyků C.

\section{IDE / Code editor}

IDE (Integrated Development Environment) je komplexní vývojové prostředí určené pro programování. Spojuje v sobě editor kódu, kompilátor nebo interpret, debugger a další nástroje potřebné pro vývoj aplikací. Umožňuje programátorovi psát, spouštět, testovat a ladit programy na jednom místě. Typickými příklady jsou například Visual Studio nebo IntelliJ IDEA.

Code editor je jednodušší nástroj sloužící především k psaní a úpravě zdrojového kódu. Obvykle nabízí zvýrazňování syntaxe, automatické doplňování kódu a možnost instalace rozšíření. Na rozdíl od IDE většinou neobsahuje všechny pokročilé nástroje pro ladění a správu projektu v základní výbavě. Mezi známé editory patří například Visual Studio Code nebo Sublime Text.

Zpočátku bylo využíváno Visual Studio 2022 při práci v systému Windows. Při přechodu na Linux byla současně provedena migrace na Visual Studio Code, do kterého byla nainstalována rozšíření pro programování v Unity. Práce ve VS Code byla shledána jako jednodušší a přehlednější, a proto nebyla provedena zpětná migrace na Visual Studio.

\section{Kolaborace}

Vzhledem k tomu, že pracujeme ve více lidech na jednom projektu, bylo potřeba vybrat způsob kolaborace. Unity nabízí svou vlastní možnost, neboli Unity Version Control. Ovšem kvůli různým problémům, které nastaly při testování této možnosti, bylo rozhodnuto o upuštění od tohoto specifického řešení. Nakonec bylo rozhodnuto o použití systémů Git a GitHub [3].

Git je bezplatný, snadno použitelný systém pro správu zdrojového kódu, který je navržen pro rychlé a efektivní zpracování menších i velmi rozsáhlých projektů. Git software disponuje nízkou technickou náročností a velmi vysokou rychlostí, čímž díky možnostem lokálního větvení a alternativních verzí zajišťuje svižný průběh jakékoli programátorské práce.

GitHub je webové rozhraní, které umožňuje více lidem současně provádět změny na jednom projektu pomocí softwaru Git. GitHub nabízí spolupráci v reálném čase a je proto nezbytným rozšířením systému Git.
